\chapter[Kriptografi: Modular sampai Kuantum]{Kriptografi: Bilangan Bulat Modular sampai Ancaman Kuantum}

Seorang mahasiswa mengumpulkan tugas berisi sandi buatan sendiri. Setiap huruf diganti nomornya,
nomor itu dikalikan tiga, lalu sisanya dibagi dua puluh enam dipakai sebagai huruf baru. Untuk
membuka pesan, ia mengalikan lagi dengan sembilan. Ketika ditanya dari mana angka sembilan itu,
jawabannya hampir benar: ``karena membagi tidak ada di sini, saya mengalikan dengan
kebalikannya.'' Penguji lalu mengganti pengali tiga dengan pengali empat. Tugas itu tidak pernah
selesai. Pada pengali empat, huruf pertama dan huruf keempat belas menghasilkan simbol yang
sama, sehingga dua pesan berbeda menghasilkan sandi identik.

Pertanyaan yang muncul dari kegagalan itu adalah pertanyaan seluruh bab ini: kapan sebuah
bilangan boleh dibagi di dalam aritmetika sisa? Jawabannya pendek: tepat ketika bilangan itu
tidak berbagi faktor dengan modulusnya. Jawaban pendek itu, ditekan terus, menghasilkan tiga hal
yang dipakai setiap peramban: pertukaran kunci, kunci publik, dan tanda tangan digital.

Tiga bab sebelumnya dipakai secara eksplisit. Bab 2 memberi bahasa untuk bijeksi dan partisi;
di sini bijeksi membuktikan bahwa perkalian dengan unit hanya mengacak urutan, dan partisi
menghitung orde. Bab 3 memberi sifat kelengkapan bilangan real, dan sifat itu tidak menyeberang,
karena $\mathbb{Z}_n$ tidak punya urutan yang cocok dengan perkalian; yang menggantikannya
adalah keterhinggaan, sebab pada himpunan berhingga pemetaan injektif otomatis surjektif. Bab 28
menjelaskan mengapa ``aman'' tidak bisa menjadi teorema tanpa syarat.

Satu peringatan berlaku untuk seluruh halaman. Semua modulus di sini kecil, dipilih agar bisa
diperiksa dengan tangan dan Python dalam hitungan detik. Modulus terbesar, $1\,050\,589$,
difaktorkan skrip enam baris. Tidak ada satu pun sandi di sini yang aman dipakai.

\begin{untukanak}
Bayangkan jam dinding di ruang kelas. Jarumnya bergerak dari angka satu sampai dua belas, lalu
kembali ke satu. Jam itu tidak pernah menunjukkan angka tiga belas atau dua puluh. Kalau kita
memutar jarum lima langkah, lalu lima langkah lagi, lalu lima langkah lagi, kita tidak perlu
tahu berapa lima ditambah lima ditambah lima; kita cukup melihat di mana jarum berhenti.

Sekarang bayangkan kunci kombinasi dengan tiga belas posisi. Kita boleh memutar sejauh berapa
pun, dan yang tercatat hanya posisi terakhir. Kalau kita memutar tiga langkah, berapa kali
kita harus memutar tiga langkah lagi supaya kembali ke tempat semula? Jawabannya sembilan, dan
itu bukan kebetulan, melainkan karena tiga dan tiga belas tidak punya faktor bersama. Coba hal
yang sama dengan empat langkah pada kunci dua belas posisi: berapa pun kali kita memutar, kita
tidak akan pernah mendarat di semua posisi, karena empat dan dua belas punya faktor bersama.

Di bab ini kita akan tahu mengapa perbedaan kecil itu menentukan apakah sebuah pesan bisa
dibuka kembali atau hilang selamanya, dan mengapa dua orang yang belum pernah bertemu bisa
menyepakati satu angka rahasia di depan orang ketiga yang mendengar semuanya.
\end{untukanak}

\section{Kelas sisa: aritmetika yang berputar}
\tingkat{dasar}

Kita mulai dengan membakukan gagasan ``sisa'' yang sudah dipakai di paragraf pertama.
\index{aritmetika modular}
Dua bilangan bulat disebut kongruen modulo $n$ bila selisihnya habis dibagi $n$.

\begin{definisi}
Misalkan $n\ge1$ bilangan bulat. Untuk $a,b\in\mathbb{Z}$ kita tulis $a\equiv b\pmod{n}$
bila $n$ membagi $a-b$. Himpunan kelas ekuivalensi di bawah relasi ini disebut
$\mathbb{Z}_n$, dan kelas yang memuat $a$ ditulis $[a]_n$.
\end{definisi}

Kita berutang dua hal: bahwa ini benar-benar relasi ekuivalensi, dan bahwa sisanya tunggal.

\begin{lemma}
Kongruensi modulo $n$ adalah relasi ekuivalensi pada $\mathbb{Z}$, dan setiap kelas
$[a]_n$ memuat tepat satu wakil pada $\{0,1,\dots,n-1\}$, sehingga $|\mathbb{Z}_n|=n$.
\end{lemma}

\begin{proof}
Refleksif: $n$ membagi $0$. Simetris: bila $n\mid(a-b)$ maka $n\mid(b-a)$. Transitif: bila
$n\mid(a-b)$ dan $n\mid(b-c)$ maka $n\mid(a-c)$ karena $(a-c)=(a-b)+(b-c)$. Untuk
ketunggalan, ambil $a$ dan tinjau $S=\{a-qn\ge0: q\in\mathbb{Z}\}$. Himpunan itu tidak kosong
menurut sifat Archimedes (Bab 3), jadi punya anggota terkecil $r$. Bila $r\ge n$ maka
$r-n\in S$ lebih kecil dari $r$, kontradiksi; jadi $0\le r<n$. Bila $r'$ juga memenuhi
$0\le r'<n$, maka $r-r'$ kelipatan $n$ yang nilainya di antara $-(n-1)$ dan $n-1$, sehingga
harus nol.
\end{proof}

Operasi pada $\mathbb{Z}_n$ didefinisikan lewat wakilnya, dan definisi semacam ini berbahaya:
bila hasilnya bergantung pada wakil, operasinya bukan fungsi. Lemma berikut menutup lubang itu
dengan satu baris aljabar.

\begin{lemma}
Bila $a\equiv a'$ dan $b\equiv b'$ modulo $n$, maka $a+b\equiv a'+b'$ dan $ab\equiv a'b'$
modulo $n$.
\end{lemma}

\begin{proof}
Tulis $a'=a+kn$ dan $b'=b+ln$. Maka $a'+b'=(a+b)+(k+l)n$ dan $a'b'=ab+(al+bk+kln)n$.
\end{proof}

Penjumlahan dan perkalian karena itu terdefinisi baik, dan keduanya mewarisi asosiatif serta
distributif dari $\mathbb{Z}$. Yang hilang adalah urutan: tidak ada cara memberi tanda
``$<$'' pada $\mathbb{Z}_n$ sehingga $a<b$ berakibat $a+c<b+c$ dan $ac<bc$ untuk $c$ positif.
Alat supremum Bab 3 berhenti di pintu masuk, dan yang bekerja di dalam adalah pencacahan.

Pembagian belum tentu ada. Kelas $[3]_{26}$ punya kebalikan, yaitu $[9]_{26}$, sebab
$3\cdot9=27\equiv1$. Kelas $[2]_{26}$ tidak punya: setiap kelipatan dua genap, sedangkan
$27$ ganjil, jadi $2b\equiv1$ tidak punya penyelesaian. Akibatnya perkalian dengan dua mod dua
puluh enam bukan bijeksi, karena $[0]$ dan $[13]$ keduanya menuju $[0]$.

\kamus{kelas sisa}{residue class}{semua bilangan yang berhenti di posisi jam yang sama,
seperti pukul tiga pagi dan pukul tiga sore}

\begin{definisi}
Kelas $[a]_n$ disebut unit bila ada $[b]_n$ dengan $[a]_n[b]_n=[1]_n$. Himpunan semua unit
modulo $n$ ditulis $U(n)$, dan banyak anggotanya ditulis $\varphi(n)$.
\end{definisi}

Untuk $n=26$ ada dua belas unit, yaitu kelas dari $1,3,5,7,9,11,15,17,19,21,23,25$. Program di
lab menghitungnya ulang; angka itu muncul lagi di bagian tiga dengan nama berbeda.

\section{Faktor persekutuan, Euklid, dan kapan pembagian ada}
\tingkat{dasar}

\index{faktor persekutuan terbesar}
Pertanyaan ``kapan $[a]_n$ punya kebalikan'' ternyata pertanyaan tentang faktor persekutuan
terbesar, yang punya algoritme berumur lebih dari dua ribu tahun.

\begin{definisi}
Untuk $a,b\in\mathbb{Z}$ yang tidak keduanya nol, $\gcd(a,b)$ adalah bilangan asli terbesar
yang membagi $a$ dan $b$.
\end{definisi}

Kunci algoritme Euklid adalah satu kesamaan himpunan: pembagi bersama $a$ dan $b$ sama persis
dengan pembagi bersama $b$ dan $a-qb$.

\begin{lemma}
Untuk setiap $q\in\mathbb{Z}$, $\gcd(a,b)=\gcd(b,a-qb)$.
\end{lemma}

\begin{proof}
Bila $d$ membagi $a$ dan $b$, maka $d$ membagi $a-qb$. Sebaliknya bila $d$ membagi $b$ dan
$a-qb$, maka $d$ membagi $(a-qb)+qb=a$. Jadi kedua pasangan punya himpunan pembagi bersama
yang sama, sehingga pembagi terbesarnya sama.
\end{proof}

Menukar $(a,b)$ menjadi $(b,a\bmod b)$ karena itu tidak mengubah jawaban tetapi mengecilkan
$b$, dan prosesnya berhenti karena bilangan asli tidak mengecil tanpa batas.

\begin{teorema}[Bezout]
Untuk $a,b$ yang tidak keduanya nol, ada $x,y\in\mathbb{Z}$ dengan $ax+by=\gcd(a,b)$.
\end{teorema}

\begin{proof}
Tinjau $S=\{ax+by>0: x,y\in\mathbb{Z}\}$. Himpunan ini tidak kosong karena memuat $|a|$ atau
$|b|$, dan setiap himpunan bagian tak kosong dari bilangan asli punya anggota terkecil; inilah
bentuk diskret dari sifat infimum di Bab 3. Misalkan $d=ax_0+by_0$ anggota terkecil itu.
Bagilah $a$ oleh $d$: $a=qd+r$ dengan $0\le r<d$. Maka $r=a(1-qx_0)+b(-qy_0)$, jadi $r\in S$
bila $r>0$. Karena $d$ terkecil, $r>0$ mustahil; jadi $r=0$ dan $d\mid a$, dan dengan cara
sama $d\mid b$. Jadi $d\le\gcd(a,b)$. Sebaliknya $\gcd(a,b)$ membagi setiap bentuk $ax+by$,
khususnya $d$, sehingga $\gcd(a,b)\le d$.
\end{proof}

\begin{teorema}
Kelas $[a]_n$ adalah unit jika dan hanya jika $\gcd(a,n)=1$.
\end{teorema}

\begin{proof}
Bila $[a]_n$ unit, ada $b$ dengan $ab\equiv1$, yaitu $ab-1=kn$. Setiap pembagi bersama $a$ dan
$n$ membagi $ab$ dan $kn$, jadi membagi $1$; maka $\gcd(a,n)=1$. Sebaliknya, bila
$\gcd(a,n)=1$, teorema Bezout memberi $ax+ny=1$. Dibaca modulo $n$, kesamaan itu berbunyi
$ax\equiv1$, jadi $[x]_n$ kebalikan $[a]_n$.
\end{proof}

Bukti itu sekaligus memberi resep: jalankan Euklid, lalu susun ulang mundur. Untuk $a=1071$ dan
$b=462$ langkahnya tiga pembagian,
\[
1071=2\cdot462+147,\qquad 462=3\cdot147+21,\qquad 147=7\cdot21+0,
\]
sehingga $\gcd=21$. Menyusun mundur memberi $21=1071\cdot(-3)+462\cdot7$, dan karena faktor
persekutuannya $21$ bukan $1$, kelas $[462]_{1071}$ bukan unit.

Berapa lama algoritme ini bekerja? Setiap dua langkah sisa setidaknya terbagi dua, jadi jumlah
pembagian sebanding dengan jumlah digit. Batas tajam diberikan Lame pada 1844: tidak lebih dari
$\log_\phi(b)+1$ pembagian dengan $\phi=(1+\sqrt5)/2$, dan yang mencapainya adalah pasangan
Fibonacci berurutan. Pasangan $(F_{47},F_{46})=(2\,971\,215\,073,\ 1\,836\,311\,903)$ butuh
$45$ pembagian, sedangkan $\log_\phi(F_{46})=44{,}33$. Untuk dua ribu pasangan acak tiga puluh
dua bit, rata-ratanya $18{,}75$ dengan maksimum $32$.

\kamus{invers modular}{modular inverse}{angka yang bila dikalikan dengan bilangan asal
menghasilkan sisa satu, seperti tombol ``batalkan'' pada mesin kasir}

\section{Fermat kecil dan Euler}
\tingkat{menengah}

Bab ini butuh satu fakta tentang pangkat, dan fakta itu lahir dari pencacahan: kalikan semua
unit dengan satu unit tetap; hasilnya daftar yang sama dengan urutan berbeda.

\begin{lemma}
Bila $\gcd(a,n)=1$, pemetaan $[x]_n\mapsto[ax]_n$ adalah bijeksi dari $U(n)$ ke $U(n)$.
\end{lemma}

\begin{proof}
Hasilnya berada di $U(n)$ karena hasilkali dua unit adalah unit: bila $xx'\equiv1$ dan
$yy'\equiv1$ maka $(xy)(x'y')\equiv1$. Pemetaannya injektif: bila $ax\equiv ay$, kalikan kedua
ruas dengan kebalikan $a$ yang ada menurut teorema unit di bagian dua. Himpunannya berhingga,
dan pada himpunan berhingga pemetaan injektif dari himpunan ke dirinya sendiri otomatis
surjektif (Bab 2). Jadi ia bijeksi.
\end{proof}

\begin{teorema}[Euler]
Bila $\gcd(a,n)=1$, maka $a^{\varphi(n)}\equiv1\pmod n$.
\end{teorema}

\begin{proof}
Misalkan $u_1,\dots,u_k$ semua unit modulo $n$, dengan $k=\varphi(n)$, dan $P=u_1\cdots u_k$.
Menurut lemma di atas, daftar $au_1,\dots,au_k$ memuat setiap unit tepat sekali, jadi
\[
a^kP\equiv(au_1)(au_2)\cdots(au_k)\equiv u_1u_2\cdots u_k=P\pmod n .
\]
Karena $P$ sendiri unit, ia punya kebalikan, dan kedua ruas boleh dikalikan dengannya.
Hasilnya $a^k\equiv1$.
\end{proof}

\begin{akibat}[Fermat kecil]
Bila $p$ prima dan $p\nmid a$, maka $a^{p-1}\equiv1\pmod p$. Untuk semua $a$ berlaku
$a^p\equiv a\pmod p$.
\end{akibat}

Bentuk kedua mengikuti yang pertama: bila $p\mid a$ kedua ruas nol, bila tidak, kalikan
$a^{p-1}\equiv1$ dengan $a$. Nama ``kecil'' membedakannya dari teorema terakhir Fermat, yang jauh lebih sulit.

Angka dua belas dari bagian pertama kini punya dua nama: $\varphi(26)=12$, dan
$3^{12}\equiv1\pmod{26}$; program memeriksa bahwa $5^{12}$ dan $7^{12}$ juga $\equiv1$.
Pangkat besar karena itu bisa dipangkas: untuk $7^{222}\bmod13$ tulis $222=12\cdot18+6$,
sehingga $7^{222}\equiv7^6\equiv9\cdot10\equiv12$ karena $7^2\equiv10$ dan $7^4\equiv9$.

Kebalikannya tidak berlaku, dan itu penting untuk kejujuran bab ini. Bilangan $561=3\cdot11\cdot17$
bukan prima, tetapi $a^{560}\equiv1\pmod{561}$ untuk setiap $a$ yang tidak berbagi faktor
dengannya; program memeriksa seluruh $a$ tanpa menemukan penyangkal. Bilangan semacam ini
disebut Carmichael, dan karena itu uji keprimaan praktis memakai Miller--Rabin.

\kamus{fungsi phi}{Euler totient function}{banyaknya posisi pada kunci kombinasi yang bisa
mencapai seluruh posisi lain, yaitu posisi yang tidak berbagi faktor dengan jumlah posisi}

\section{Grup siklik, orde, dan akar primitif}
\tingkat{menengah}

Kalau kita mengalikan satu unit berulang kali, hasilnya berputar dan akhirnya kembali ke satu.
Pertanyaan yang berguna adalah berapa lama putaran itu.

\begin{definisi}
Untuk $\gcd(a,n)=1$, orde $a$ modulo $n$ adalah bilangan asli terkecil $k$ dengan
$a^k\equiv1\pmod n$, ditulis $\operatorname{ord}_n(a)$.
\end{definisi}

Orde selalu ada: barisan $a^0,a^1,\dots$ hanya punya $\varphi(n)$ nilai berbeda, jadi dua di
antaranya sama dan $a^{i-j}\equiv1$ untuk suatu $i>j$. Batas itu bisa dipertajam: orde membagi
$\varphi(n)$.

\begin{lemma}
Bila $\gcd(a,n)=1$, maka $\operatorname{ord}_n(a)$ membagi $\varphi(n)$.
\end{lemma}

\begin{proof}
Misalkan $k=\operatorname{ord}_n(a)$ dan $H=\{a^0,\dots,a^{k-1}\}$, yang anggotanya berbeda
semua, jadi $|H|=k$. Untuk setiap unit $x$, himpunan $xH=\{xa^0,\dots,xa^{k-1}\}$ juga punya
$k$ anggota, karena perkalian dengan unit adalah bijeksi. Himpunan-himpunan $xH$ membentuk
partisi dari $U(n)$: dua di antaranya sama atau lepas, sebab $xH$ dan $yH$ beririsan bila dan
hanya bila $y^{-1}x\in H$ (Bab 2). Jadi $\varphi(n)$ adalah $k$ dikalikan banyaknya bagian.
\end{proof}

\begin{definisi}
Bila $\operatorname{ord}_n(g)=\varphi(n)$, maka $g$ disebut akar primitif modulo $n$. Untuk
unit $h$, bilangan $k$ dengan $g^k\equiv h$ disebut logaritma diskret $h$ terhadap basis $g$.
\end{definisi}

Modulo $13$, $2,6,7,11$ berorde penuh $12$ dan karena itu akar primitif, sedangkan $3$ dan $9$
berorde $3$, $5$ dan $8$ berorde $4$, $4$ dan $10$ berorde $6$, dan $12$ berorde $2$. Semua
angka itu membagi $12$.

Bagi $2$, putarannya mengunjungi keduabelas kelas bukan nol sebelum kembali:
$2,4,8,3,6,12$, lalu $11,9,5,10,7,1$. Gambar~\ref{fig:jam} menampilkannya sebagai jam tiga
belas posisi. Bagi $3$, putarannya hanya $3,9,1$ lalu berulang.

\begin{figure}[htbp]
\centering
\begin{tikzpicture}
\foreach \k in {0,...,12}
  \node[circle,draw,inner sep=1pt,minimum size=6mm,font=\scriptsize] (n\k)
    at ({90-360*\k/13}:2.1cm) {\k};
\node[font=\footnotesize,align=center] at (0,0) {kalikan\\ dengan 2};
\foreach \a/\b in {1/2,2/4,4/8,8/3,3/6,6/12,12/11,11/9,9/5,5/10,10/7,7/1}
  \draw[-{Latex[length=1.5mm]},semithick] (n\a) to[bend left=18] (n\b);
\end{tikzpicture}
\caption{Perkalian dengan $2$ modulo $13$ berputar melalui keduabelas kelas bukan nol
sebelum kembali ke $1$. Panjang putaran itulah orde, dan orde penuh berarti $2$ akar
primitif modulo $13$.}
\label{fig:jam}
\end{figure}

Keberadaan akar primitif untuk setiap prima $p$ adalah teorema Gauss, dan buktinya kami
lewati. Yang perlu dibawa: untuk prima $p$, unit adalah satu putaran berisi $p-1$ posisi, dan
setiap unit adalah pangkat dari satu akar primitif. Bila $h=g^k$, mengetahui $k$ berarti
mengetahui $h$, tetapi tidak sebaliknya; asimetri itu bahan mentah kriptografi kunci publik.

\kamus{akar primitif}{primitive root}{satu posisi awal yang, bila diputar terus, melewati
semua posisi lain sebelum kembali}

\section{Pertukaran kunci Diffie-Hellman}
\tingkat{menengah}

\index{diffie-hellman}
Dua orang ingin menyepakati satu angka rahasia, dan satu-satunya saluran yang mereka punya
didengar pihak ketiga. Pada 1976 Diffie dan Hellman menunjukkan bahwa hal itu mungkin, dan
mekanismenya hanya memakai dua rumus.

Mereka sepakat lebih dulu pada prima $p$ dan akar primitif $g$ modulo $p$; keduanya boleh
diumumkan. Alice memilih $a$ rahasia dan mengirim $A=g^a\bmod p$. Bob memilih $b$ rahasia dan
mengirim $B=g^b\bmod p$. Alice menghitung $K=B^a\bmod p$, Bob menghitung $K=A^b\bmod p$.
Keduanya memperoleh angka yang sama, karena
\[
B^a\equiv(g^b)^a\equiv g^{ab}\equiv(g^a)^b\equiv A^b\pmod p .
\]
Yang dilihat penyadap hanyalah $p,g,A,B$. Untuk memperoleh $K$ ia perlu $g^{ab}$, dan cara
termudah yang diketahui adalah menemukan $a$ atau $b$, yaitu menyelesaikan logaritma diskret.

\begin{bacakeras}
Bacakan sebagai dialog dua suara, bergantian, satu baris satu tarikan napas.

Suara pertama: ``Sepakat dulu, umumkan dua.''
Suara kedua: ``Simpan sendiri, jangan disebut.''
Suara pertama: ``Pangkatkan, bagi sisa, kirim hasilnya.''
Suara kedua: ``Terima kiriman, pangkatkan lagi.''
Suara pertama: ``Sisa terakhir, itulah kuncinya.''
Bersama: ``Dua arah, satu angka, tidak pernah lewat.''

Ulangi tiga kali, makin cepat. Baris yang membuat lidah tersendat adalah langkah yang belum
dipahami, bukan baris yang perlu dihapal lebih keras. Sesudah lancar, tutup buku dan sebutkan
empat hal berurutan: yang diumumkan, yang disimpan, yang dikirim, yang dihitung. Menyebut angka
rahasia lebih dulu membatalkan seluruh permainan.

Baris yang diulang tiga kali dengan suara keras: ``yang dikirim hasil pangkat, yang disimpan
pangkatnya.''
\end{bacakeras}

Angka konkretnya, seluruhnya diperiksa program di bagian lab: $p=10\,007$, $g=5$ berorde
$10\,006$ sehingga akar primitif, $a=1234$, $b=5678$. Maka $A=5^{1234}\bmod p=4873$ dan
$B=5^{5678}\bmod p=3189$. Alice menghitung $3189^{1234}\bmod p=5674$; Bob menghitung
$4873^{5678}\bmod p=5674$. Kunci bersama mereka $5674$, dan angka itu tidak pernah melintasi
saluran. Gambar~\ref{fig:dh} menampilkan alurnya beserta apa yang didengar penyadap.

\begin{figure}[htbp]
\centering
\begin{tikzpicture}[font=\footnotesize,
  kotak/.style={draw,rounded corners=2pt,align=left,inner sep=4pt,text width=4.4cm}]
\node[kotak] (a) {\textbf{Alice}\\ $a=1234$ rahasia\\ $A=5^{1234}\bmod p=4873$};
\node[kotak,right=3.4cm of a] (b) {\textbf{Bob}\\ $b=5678$ rahasia\\
  $B=5^{5678}\bmod p=3189$};
\node[kotak,below=1.3cm of $(a)!0.5!(b)$, text width=9.2cm, fill=black!4]
  (e) {\textbf{Penyadap} mendengar $p=10\,007$, $g=5$, $A=4873$, $B=3189$, lalu bermaksud
  menghitung $K=A^{b}$ atau $K=B^{a}$; untuk itu ia harus menemukan $a$ atau $b$.};
\draw[-{Latex}] ([yshift=3mm]a.east) -- node[above,font=\scriptsize]{$A$}
  ([yshift=3mm]b.west);
\draw[-{Latex}] ([yshift=-3mm]b.west) -- node[below,font=\scriptsize]{$B$}
  ([yshift=-3mm]a.east);
\draw[-{Latex},dashed] (e.north -| a.south) -- (a.south);
\draw[-{Latex},dashed] (e.north -| b.south) -- (b.south);
\end{tikzpicture}
\caption{Alur pertukaran kunci Diffie-Hellman pada contoh angka di teks. Nilai $A$ dan $B$
melintasi saluran terbuka; yang tidak melintas adalah eksponen $a$ dan $b$, dan justru
keduanya yang dibutuhkan penyadap.}
\label{fig:dh}
\end{figure}

Kesamaan itu hanya menyatakan protokolnya konsisten, bukan aman. Yang bisa dibuktikan hanyalah
pernyataan bersyarat: bila logaritma diskret modulo $p$ sulit, menghitung $g^{ab}$ dari $A$
dan $B$ juga sulit. Syarat itu bukan teorema; ada grup tempat logaritma diskret mudah, jadi
keamanannya sifat kelompok tertentu.

\begin{tangan}
Bahan: cat air kuning, satu warna rahasia per orang (merah dan biru), empat gelas plastik,
satu sendok, dua lembar kertas putih.

Langkah pertama, di depan meja: tuang satu sendok cat kuning ke gelas masing-masing. Kuning itu
peran prima dan basis, dan keduanya diumumkan. Langkah kedua, di balik punggung: tiap orang
menambahkan dua sendok warna rahasianya, mengaduk, lalu menyerahkan gelas campurannya. Langkah
ketiga: masing-masing menambahkan dua sendok warna rahasianya sendiri ke gelas yang baru
diterima, mengaduk, lalu menotolkan satu tetes ke kertas.

Yang harus terlihat: dua totolan itu sewarna, padahal tidak satu pun dari kedua orang itu
pernah menyerahkan warna rahasianya. Gelas campuran pertama yang dipegang orang ketiga tidak
bisa dipisah kembali; ia hanya melihat dua warna kusam yang berbeda.
\end{tangan}

Dua jebakan parameter layak disebut. Pertama, $g$ harus membangkitkan subkelompok besar berorde
prima: modulus $10\,007$ adalah prima aman dengan $p-1=2\cdot5003$, dan program mencatat bahwa
$5^{5003}\bmod p=10\,006$ berorde hanya $2$. Bila Alice memakai basis itu, kunci bersamanya
hanya bernilai $1$ atau $10\,006$, dan penyadap tidak perlu menghitung logaritma apa pun.
Kedua, protokol ini tidak menyebut identitas siapa pun, sehingga penyadap bisa berdiri di tengah
dan membaca semuanya. Gambar~\ref{fig:dh} hanya menggambarkan penyadap pasif.

Tentang ukuran: modulus $10\,007$ bisa dibobol dalam waktu kurang dari satu detik oleh
algoritme indeks kalkulus, dan kita memakainya karena hasil pangkatnya muat di satu halaman.
Praktik nyata memakai $2048$ bit atau lebih, dan itu pun batas waktu, bukan jaminan.

\begin{untukanak}
Bayangkan dua orang yang ingin menyepakati satu warna rahasia, padahal semua percakapan mereka
didengar. Mereka sepakat lebih dulu pada satu warna dasar, misalnya kuning. Alice memilih warna
rahasianya, mencampurnya dengan kuning, lalu mengirim hasil campurannya. Bob melakukan hal yang
sama dengan warna rahasianya sendiri. Sekarang masing-masing mencampur warna yang diterimanya
dengan warna rahasianya sendiri. Keduanya mendapat warna yang sama, dan penyadap hanya melihat
dua campuran yang tidak bisa dipisah kembali.

Pemisahan itulah yang dianggap sulit: mencampur warna gampang, memisahkan satu campuran tanpa
salah satu asalnya dianggap tidak. Kriptografi mengganti warna menjadi angka dan ``dianggap
tidak gampang'' menjadi kalimat matematika.

Di bab ini kita akan tahu bagaimana angka menggantikan warna, dan mengapa keamanannya
bergantung pada satu kalimat ``dianggap'' yang tidak pernah dibuktikan siapa pun.
\end{untukanak}

\section{RSA: dari kunci ke pesan dan kembali}
\tingkat{menengah}

\index{rsa}
Diffie-Hellman menyepakati kunci, tetapi tidak mengirim pesan. Setahun kemudian Rivest,
Shamir, dan Adleman menunjukkan bahwa sifat yang sama bisa dipakai untuk dua kunci berbeda:
satu untuk mengunci, satu untuk membuka.

Pilih dua prima besar $p$ dan $q$, hitung $n=pq$ dan $\varphi(n)=(p-1)(q-1)$. Pilih $e$ dengan
$\gcd(e,\varphi(n))=1$; teorema unit memberi kebalikannya $d$ modulo $\varphi(n)$. Kunci publik
$(n,e)$, kunci rahasia $d$. Untuk pesan $0\le m<n$, sandinya $c=m^e\bmod n$ dan pembukaannya
$m=c^d\bmod n$.

Mengapa pembukaan mengembalikan $m$? Karena $ed=1+k\varphi(n)$, sehingga
$c^d=m^{ed}=m\,(m^{\varphi(n)})^k$. Bila $\gcd(m,n)=1$, teorema Euler membuat kurung itu satu.
Kasus sisanya menjelaskan mengapa $n$ harus hasilkali dua prima berbeda.

\begin{teorema}
Misalkan $n=pq$ dengan $p,q$ prima berbeda, $ed\equiv1\pmod{(p-1)(q-1)}$, dan $0\le m<n$.
Maka $m^{ed}\equiv m\pmod n$.
\end{teorema}

\begin{proof}
Bila $p\nmid m$, Fermat kecil memberi $m^{p-1}\equiv1\pmod p$, sehingga
$m^{ed}=m\,(m^{p-1})^{k(q-1)}\equiv m\pmod p$. Bila $p\mid m$, kedua ruas kongruen nol modulo
$p$. Jadi kesamaan berlaku modulo $p$ untuk setiap $m$, dan dengan alasan sama modulo $q$.
Karena $p$ dan $q$ keduanya membagi $m^{ed}-m$ dan $\gcd(p,q)=1$, hasilkali keduanya juga
membagi bilangan itu: bila $m^{ed}-m=pt$ dan $q\mid pt$ dengan $q$ prima serta $q\nmid p$,
maka $q\mid t$, jadi $pq\mid m^{ed}-m$.
\end{proof}

Contoh angkanya kecil dan seluruhnya diperiksa program. Ambil $p=1019$, $q=1031$, maka
$n=1\,050\,589$ dan $\varphi(n)=1018\cdot1030=1\,048\,540$. Pilih $e=65\,537$; Euklid memberi
$d=60\,493$, dan memang $65\,537\cdot60\,493=3\,964\,529\,741=1+3781\cdot1\,048\,540$. Pesan
$m=424\,242$ menghasilkan sandi $c=552\,241$, dan $c^d\bmod n$ kembali $424\,242$. Modulus itu hanya dua puluh satu bit dan difaktorkan skrip enam baris.

RSA polos, rumus di atas tanpa tambahan apa pun, tidak aman; tiga alasannya bisa ditunjukkan
dengan modulus yang sama. Pertama, ia deterministik: pesan yang sama selalu menghasilkan sandi
yang sama, sehingga pesan berulang bisa dikenali. Kedua, ia multiplikatif:
$E(m_1)E(m_2)\equiv E(m_1m_2)$, dan program memperlihatkan
$E(1234)\cdot E(567)\bmod n=551\,173$ dibuka menjadi $699\,678$, tepat $1234\cdot567$. Ketiga,
pesan kecil tidak terlindungi: dengan $e=3$ dan $m=42$, sandinya $74\,088$, kebetulan sama
dengan $42^3$ karena hasilnya belum melewati $n$, sehingga akar pangkat tiga biasa memulihkan
pesan.

Ketiganya ditutup dengan bantalan: pesan diberi tambahan acak sebelum dipangkatkan, sehingga
hasilnya panjang penuh, selalu berbeda, dan tidak lagi multiplikatif. RSA-OAEP, standar yang
dipakai hari ini, punya bukti keamanan terhadap serangan chosen-ciphertext dalam model oracle
acak dengan asumsi RSA berlaku. Dalam praktik RSA membungkus kunci simetris $256$ bit lalu data
dienkripsi dengan AES; pola ini disebut enkripsi hibrida, karena RSA lambat: waktunya
sebanding pangkat tiga panjang modulus.

\kamus{bantalan}{padding}{isi tambahan yang dimasukkan ke dalam paket sebelum dikirim,
seperti bubble wrap yang membuat setiap paket berbeda bentuk dan ukuran}

\section{Fungsi hash, MAC, dan tanda tangan}
\tingkat{menengah}

\index{fungsi hash}
Enkripsi menyembunyikan isi. Kadang yang dibutuhkan bukan menyembunyikan, melainkan meringkas:
satu sidik jari pendek yang berubah bila isinya berubah.

\begin{definisi}
Fungsi hash kriptografis memetakan pesan sepanjang apa pun ke keluaran berukuran tetap. Ia
harus tahan pra-gambar (dari $h$, sulit menemukan $m$ dengan $H(m)=h$), tahan pra-gambar
kedua (dari $m$, sulit menemukan $m'\neq m$ dengan $H(m')=H(m)$), dan tahan tabrakan (sulit
menemukan sepasang $m,m'$ berbeda dengan hash sama).
\end{definisi}

Ketiga sifat itu tidak setara; yang paling lemah paling sering disalahpahami. Ketahanan
tabrakan dibatasi ulang tahun: bila keluaran punya $N$ nilai, tabrakan muncul setelah kira-kira
$1{,}177\sqrt N$ percobaan, bukan setelah $N$; untuk $128$ bit angka itu $2{,}17\cdot10^{19}$
dan untuk $64$ bit $5{,}06\cdot10^{9}$. Program di lab menguji ramalan itu pada SHA-256 yang
dipotong menjadi $16$ bit: rata-rata $362{,}7$ sampel melawan ramalan $320{,}8$, dengan sebaran
dari $100$ sampai $743$. Jadi hash berkeluaran $n$ bit hanya memberi ketahanan tabrakan $n/2$
bit.

MAC menambahkan kunci pada hash: dua pihak yang berbagi kunci dapat memeriksa bahwa pesan
datang dari pihak lain dan tidak berubah, karena tanpa kunci tidak ada cara menghitung tag yang
benar. Batasnya: MAC hanya berguna antara dua pihak yang sama-sama memegang kunci, sehingga
pihak ketiga tidak bisa menjadi saksi.

Tanda tangan memakai pasangan kunci dengan arah terbalik: penandatangan memakai kunci rahasia,
pemeriksa memakai kunci publik. Pada RSA, tanda tangan atas $m$ adalah $s=H(m)^d\bmod n$, dan
pemeriksaan membandingkan $s^e\bmod n$ dengan $H(m)$. Keharusan memakai $H$ bukan soal
kenyamanan: tanpa hash, tanda tangan RSA bersifat multiplikatif seperti enkripsinya,
sehingga dari tanda tangan dua dokumen seseorang bisa menyusun tanda tangan sah untuk hasilkali
keduanya tanpa memiliki kunci rahasia. Hash memutus sifat itu, karena $H(m_1m_2)$ tidak
berhubungan dengan $H(m_1)H(m_2)$. Syarat keamanannya adalah ketahanan terhadap pemalsuan
eksistensial di bawah serangan pesan terpilih, dan sekali lagi itu asumsi.

\kamus{tanda tangan digital}{digital signature}{cap lilin pada dokumen yang bisa diperiksa
siapa saja tetapi hanya bisa dibuat oleh pemegang cetakannya}

\section{Apa yang sebenarnya diasumsikan model keamanan}
\tingkat{lanjut}

Setiap kalimat ``sulit'' di bab ini adalah janji tentang apa yang tidak bisa dilakukan
algoritme yang dikenal, bukan pernyataan tentang apa yang mustahil.

\begin{definisi}
Fungsi $f$ disebut satu arah bila ada algoritme polinomial untuk menghitung $f(x)$ dari $x$,
tetapi untuk $y$ acak di daerah hasil, setiap algoritme polinomial gagal menemukan $x$ dengan
$f(x)=y$ kecuali dengan peluang yang dapat diabaikan.
\end{definisi}

Arahnya digambarkan pada Gambar~\ref{fig:tangga-satu-arah}.

\begin{figure}[htbp]
\centering
\begin{tikzpicture}[font=\small,
  kotak/.style={draw,rounded corners=2pt,align=center,inner sep=5pt,
    minimum height=0.95cm}]
  \node[kotak,draw=layerbrown] (pq) at (0,0) {$p$, $q$ prima acak};
  \node[kotak,draw=ink] (n) at (4.0,0) {$n=p\,q$};
  \node[kotak,draw=layergrey,dashed] (s) at (8.2,0) {penyerang};
  \draw[-{Latex},very thick,layerjade] (pq) -- (n);
  \node[font=\scriptsize,layerjade,above] at (2.0,0.06)
    {kalikan: mudah, polinomial};
  \draw[-{Latex},thick,layergrey] (n) -- (s);
  \node[font=\scriptsize,layergrey,above] at (6.1,0.06) {melihat $n$};
  \draw[layerplum,thick,dashed] (n.south) -- (4.0,-1.35);
  \draw[layerplum,thick,dashed] (0,-1.35) -- (pq.south);
  \draw[-{Latex},layerplum,thick,dashed] (4.0,-1.35) -- (0,-1.35);
  \draw[layerplum,very thick] (1.86,-1.12)--(1.86,-1.58);
  \draw[layerplum,very thick] (2.04,-1.12)--(2.04,-1.58);
  \node[font=\scriptsize,layerplum,below] at (2.0,-1.62)
    {faktorkan: sulit, tidak ada algoritme polinomial yang dikenal};
\end{tikzpicture}
\caption{Dua arah yang tidak sama mahal. Dari $p$ dan $q$ ke $n$ cukup satu
perkalian, sedangkan dari $n$ kembali ke $p$ dan $q$ adalah masalah yang sampai
hari ini tidak punya algoritme polinomial yang dikenal; selisih itulah yang
dipakai RSA dan Diffie-Hellman. Palang pada panah balik menandai arah yang
diandaikan tertutup, bukan arah yang sudah terbukti mustahil.}
\label{fig:tangga-satu-arah}
\end{figure}

Definisi itu tidak menyebut satu fungsi pun, dan itu bukan kelalaian. Sampai hari ini tidak ada
bukti bahwa fungsi satu arah ada; keberadaannya akan berakibat $P\neq NP$, yang juga masih
terbuka. Kriptografi kunci publik karena itu berdiri di atas konjektur, dan rantainya berakhir di
asumsi yang tidak berubah menjadi teorema karena dipakai lama.

\index{asumsi faktorisasi}
Asumsi faktorisasi mengatakan: untuk $n$ hasilkali dua prima acak berukuran $k$ bit, tidak ada
algoritme polinomial dalam $k$ untuk memfaktorkannya. Algoritme terbaik yang diketahui,
penyaringan medan bilangan, berjalan dalam waktu $L_n[1/3,1{,}923]$, yang tumbuh
subeksponensial. Bilangan RSA-250, sepanjang $250$ digit atau sekitar $829$ bit, difaktorkan
pada 2020 dengan biaya sekitar $2\,700$ core-tahun.

Padanan kekuatan yang dicatat NIST dalam SP 800-57 berguna sebagai peta kasar.

\begin{center}
\small
\begin{tabular}{@{}lll@{}}
\toprule
Kekuatan (bit) & Modulus RSA atau DH (bit) & Sandi simetris \\
\midrule
112 & 2048 & 3DES tiga kunci \\
128 & 3072 & AES-128 \\
192 & 7680 & AES-192 \\
256 & 15360 & AES-256 \\
\bottomrule
\end{tabular}
\end{center}

Angka kolom tengah adalah taksiran biaya algoritme terbaik saat tabel disusun, jadi tabel
semacam ini punya tanggal kedaluwarsa.

Model keamanan formal menutup celah retorika dengan permainan: penantang membangkitkan kunci,
menjawab permintaan enkripsi, lalu memilih dua pesan dan mengirimkan enkripsi salah satunya;
skema aman bila penyerang tidak bisa menebak yang mana jauh di atas peluang setengah. Keamanan
dibuktikan sebagai reduksi: penyerang yang menang diubah menjadi algoritme pemfaktor $n$,
sehingga seluruh keraguan pindah ke satu tempat yang sempit tetapi tidak hilang.

\section{Shor, dan apa yang runtuh bersamanya}
\tingkat{lanjut}

\index{algoritma shor}
Seluruh bangunan di atas bergantung pada satu kesulitan, dan pada 1994 Peter Shor menunjukkan
bahwa komputer kuantum menyelesaikannya dalam waktu polinomial.

Gagasannya terbagi dua, dan bagian pertama klasik sepenuhnya. Untuk memfaktorkan $n$, pilih $a$
dengan $\gcd(a,n)=1$ dan hitung $r=\operatorname{ord}_n(a)$. Bila $r$ ganjil, ulangi dengan $a$
lain. Bila $r$ genap dan $a^{r/2}\not\equiv-1\pmod n$, maka $\gcd(a^{r/2}-1,n)$ adalah faktor
tak sejati dari $n$: karena $a^r\equiv1$, bilangan $a^{r/2}$ adalah akar kuadrat dari satu
yang bukan $\pm1$, dan akar semacam itu hanya muncul bila modulusnya terbagi. Sisa pekerjaannya
mencari orde, dan inilah bagian yang kuantum.

Bagian klasik saja tidak cukup. Untuk $n=1019\cdot1031$ kami mencoba $2\,996$ unit acak dan
hanya $319$, sekitar $10{,}6$ persen, memenuhi syarat sehingga memberi faktor; keberhasilan
karena itu hasil pengulangan dengan $a$ berbeda. Yang hilang tanpa komputer kuantum adalah cara
cepat menemukan $r$: tidak ada algoritme klasik yang diketahui berjalan polinomial dalam
$\log n$.

Komputer kuantum mencari $r$ dengan transformasi Fourier kuantum, dan di situlah percepatannya
berasal. Akibatnya bukan hanya RSA yang runtuh: algoritme Shor juga menyelesaikan logaritma
diskret, baik pada $\mathbb{Z}_p$ maupun pada kurva eliptik, sehingga Diffie-Hellman, DSA, ECDH, dan ECDSA ikut runtuh, termasuk kurva eliptik yang modulusnya paling
pendek karena kebal terhadap indeks kalkulus.

Yang tidak runtuh adalah sandi simetris dan fungsi hash. Grover memberi percepatan kuadratik
pada pencarian kunci, sehingga AES-128 turun dari $2^{128}$ menjadi sekitar $2^{64}$ operasi
kuantum, sedangkan AES-256 tetap di sekitar $2^{128}$; ketahanan tabrakan hash praktis tidak
berubah. Jawaban untuk data simetris karena itu sederhana: pakai kunci dua kali lebih panjang.

Perlu dibedakan dua hal yang sering dicampur dalam berita. Shor membuktikan bahwa faktornya
bisa dihitung polinomial pada komputer kuantum tanpa galat; membangun mesinnya persoalan
teknik lain. Taksiran Gidney dan Ekera (2021) menyebut sekitar dua puluh juta qubit fisik dan
delapan jam untuk modulus $2048$ bit, hampir seluruhnya untuk koreksi galat. Angka itu turun
beberapa orde besaran dalam satu dekade, tetapi belum ada mesin yang mendekati. Yang berlaku
sekarang adalah pola ``tangkap sekarang, buka nanti'': pesan hari ini bisa disimpan dan dibuka
ketika mesinnya ada.

\section{Kriptografi pasca-kuantum: lattice dan Kyber}
\tingkat{lanjut}

\index{kriptografi pasca-kuantum}
Bila faktorisasi jatuh, yang dibutuhkan adalah masalah lain yang tidak jatuh bersamanya.
Kandidat terkuat berasal dari lattice, yaitu himpunan titik yang dibangkitkan kombinasi bilangan
bulat dari beberapa vektor.

\index{lattice}
\emph{Learning with errors} menyusun masalahnya begini. Diberikan matriks $A$ acak
$m\times n$ atas $\mathbb{Z}_q$ dan vektor $b=As+e\bmod q$, dengan $s$ rahasia dan $e$ galat
kecil, temukan $s$. Tanpa $e$ ini hanya sistem persamaan linear yang diselesaikan eliminasi
Gauss (Bab 8); dengan $e$ kecil tetapi tidak nol, sistemnya hampir konsisten dan mencari $s$
dianggap sulit. Reduksi Regev (2005) menambah daya tariknya: menyerang skema acak sama sulitnya
dengan menyelesaikan kasus terburuk lattice, bukan kasus rata-rata pilihan penyerang.

Kyber, yang dibakukan NIST pada Agustus 2024 sebagai ML-KEM dalam FIPS 203, adalah varian modul
dari gagasan itu dengan $n=256$, $k=3$, dan $q=3329$. Kunci publiknya $1\,184$ byte dan
ciphertext-nya $1\,088$ byte, sekitar $37$ dan $34$ kali kunci simetris $32$ byte yang
dibungkusnya, dan biaya ruang itu nyata. Tanda tangan pasca-kuantum punya cerita serupa:
ML-DSA (Dilithium) juga berbasis lattice, sedangkan SLH-DSA (SPHINCS+) memilih jalan
konservatif berbasis hash.

Kejujuran yang wajib di sini: asumsi lattice jauh lebih muda daripada asumsi faktorisasi.
Faktorisasi diserang sejak Gauss; LWE baru berumur dua puluh tahun, dan sejarah kriptografi
penuh skema yang tampak kuat lalu patah setelah sepuluh tahun dianalisis. Karena itu transisi
yang dianjurkan bukan penggantian, melainkan hibrida: kunci sesi dibentuk dari Kyber dan dari
pertukaran Diffie-Hellman sekaligus.

\section{Ketika keluaran menjadi oracle: padding dan model}
\tingkat{lanjut}

\index{padding oracle}
Ada satu keluarga serangan yang tidak menyerang matematikanya, melainkan cara program menjawab.
Bleichenbacher pada 1998 menunjukkan bahwa RSA dengan bantalan PKCS\#1 v1.5 bisa dibuka
sepenuhnya bila server menjawab berbeda untuk ``bantalan salah'' dan ``bantalan benar tetapi isi
salah''. Setiap jawaban membocorkan satu bit tentang pesan, dan dengan beberapa ribu pertanyaan
pesannya terbaca. Matematika RSA tidak dilanggar; yang dilanggar adalah asumsi bahwa penyerang
hanya melihat ciphertext.

Pola ancamannya sama dengan yang dibahas Bab 29 dan Bab 30. Model pembelajaran mesin yang
mengembalikan probabilitas, atau pesan galat yang berbeda-beda, adalah oracle; penyerang yang
boleh bertanya berkali-kali memakai perbedaan kecil itu untuk menyimpulkan apa yang seharusnya
tidak ia ketahui, termasuk apakah sebuah contoh ada di data latih. Bab 30 menambahkan
peringatan yang berlaku di kedua dunia: pengukuran yang tampak membaik belum tentu karena
pertahanannya bekerja.

Pelajaran defensifnya karena itu sama. Keluaran harus dibuat seragam: satu pesan galat untuk
semua kegagalan, waktu tanggap yang tidak bergantung pada rahasia, penolakan yang tidak
membedakan sebab. Untuk model, ini berarti membatasi jumlah kueri dan menyamarkan skor. Semua
itu mengurangi kebocoran, tidak menghapusnya.

\begin{tantangan}
Sebuah pustaka memakai AES-GCM dengan kunci $K$ yang sama untuk dua pesan, tetapi penghitung
nonce-nya ter-reset sehingga nonce $96$ bit yang sama terpakai dua kali. Penyerang memegang
kedua sandi, kedua tag, dan isi pesan pertama, yaitu halaman login yang tetap. Tunjukkan
bagaimana ia memperoleh kunci autentikasi $H$, lalu mengapa $H$ cukup untuk memalsukan tag pesan
buatannya tanpa mengetahui $K$.

Petunjuk 1: nonce yang sama memberi kunci aliran yang sama, sehingga selisih kedua sandi
menghapus kunci itu.

Petunjuk 2: tag adalah polinomial dalam $H$ dengan koefisien yang diketahui; dua tag dengan
nonce sama memberi dua persamaan bagi satu $H$, dan selisihnya menjadi polinomial yang akarnya
dapat dihitung.

Kalau mentok, kembali ke bagian ``Fungsi hash, MAC, dan tanda tangan'' dan baca ulang paragraf
tentang MAC.
\end{tantangan}

\section{Contoh hitung}
\tingkat{menengah}

Kita hitung satu invers dengan tangan, karena langkah mundurnya adalah tempat kesalahan paling
sering terjadi. Ambil $\varphi(n)=1\,048\,540$ dari modulus mainan di bagian enam dan
$e=65\,537$. Euklid maju memberi
\[
1\,048\,540=15\cdot65\,537+65\,485,\qquad 65\,537=1\cdot65\,485+52,
\]
\[
65\,485=1259\cdot52+17,\qquad 52=3\cdot17+1,\qquad 17=17\cdot1+0 .
\]
Sisa terakhir yang bukan nol adalah $1$, jadi $\gcd=1$ dan inversnya ada. Susun mundur mulai
dari sisa satu:
\[
1=52-3\cdot17=52-3(65\,485-1259\cdot52)=3778\cdot52-3\cdot65\,485 .
\]
Lanjutkan dengan $52=65\,537-65\,485$ dan $65\,485=1\,048\,540-15\cdot65\,537$:
\[
1=3778\cdot65\,537-3781\cdot65\,485
=3778\cdot65\,537-3781(1\,048\,540-15\cdot65\,537).
\]
Kelompokkan: $1=(3778+15\cdot3781)\cdot65\,537-3781\cdot1\,048\,540
=60\,493\cdot65\,537-3781\cdot1\,048\,540$. Jadi $d=60\,493$, dan tanda negatifnya hilang
karena dibaca modulo $\varphi(n)$. Pemeriksaannya satu baris: $65\,537\cdot60\,493
=3\,964\,529\,741$, sedangkan $3781\cdot1\,048\,540=3\,964\,529\,740$; selisihnya satu.

Dua hitungan lain di bab ini bisa diperiksa dengan cara yang sama: pemangkasan $7^{222}\bmod13$
di bagian tiga, yang memakai $222=12\cdot18+6$ dan memberi $12$, dan pertukaran kunci kecil
$p=23$, $g=5$, $a=4$, $b=3$ yang memberi $A=4$, $B=10$, dan kunci $18$ dari kedua arah.

\begin{resep}
Empat langkah mencari invers $a$ modulo $n$.
\begin{enumerate}
  \item Tulis $(a,n)$. Jalankan Euklid: ganti $(a,n)$ dengan $(n,a\bmod n)$ sampai sisanya
        nol; catat tiap baris $n_0=q\,n_1+r$.
  \item Sisa bukan nol terakhir adalah $\gcd$. Bila bukan $1$, berhenti; inversnya tidak ada.
  \item Dari baris terakhir yang bersisa $1$, tulis $1=r-q\,s$, lalu ganti $s$ dan $r$ dengan
        bentuk baris di atasnya, satu baris ke atas untuk setiap langkah.
  \item Kumpulkan koefisien $a$; itulah inversnya. Bila negatif, tambahkan $n$ sampai masuk
        $0,\dots,n-1$, lalu periksa $a\cdot x\bmod n=1$.
\end{enumerate}
Peringatan urutan: susun mundur dari sisa terkecil naik ke atas, jangan dari baris pertama.
Menukar peran $a$ dan $n$ juga mengubah jawaban: invers $a$ modulo $n$ bukan invers $n$
modulo $a$.
\end{resep}

\kertas{Ulangi perhitungan pertama tanpa melihat teks. Tulis lima baris Euklid maju untuk
$65\,537$ dan $1\,048\,540$, lalu susun mundur sampai memperoleh $d$. Periksa dengan
mengalikan $e$ dan $d$ lalu membaginya dengan $\varphi(n)$: sisanya harus satu. Sesudah itu
jawab: mengapa langkah yang sama tidak bisa dipakai untuk membuka RSA tanpa kunci rahasia,
padahal seluruh perhitungan ini hanya butuh lima pembagian?}

\section{Lab komputasi}
\tingkat{lanjut}

Program berikut memakai pustaka standar saja dan berjalan di bawah satu detik. Ia memeriksa
empat klaim di badan bab: unit modulo $26$, statistik langkah Euklid, kesamaan kunci
Diffie-Hellman, dan kelonggaran RSA polos. Bagian terakhir menguji ramalan ulang tahun pada
SHA-256 yang dipotong menjadi dua byte.

\begin{lab}
\begin{verbatim}
import hashlib, math, random

def egcd(a, b):
    if b == 0:
        return a, 1, 0
    g, x, y = egcd(b, a % b)
    return g, y, x - (a // b) * y

def invers(a, n):
    g, x, _ = egcd(a, n)
    return None if g != 1 else x % n

def langkah(a, b):
    k = 0
    while b:
        a, b = b, a % b
        k += 1
    return k

def orde(a, n):
    k, y = 1, a % n
    while y != 1:
        y = y * a % n
        k += 1
    return k

unit = [a for a in range(26) if invers(a, 26) is not None]
print("unit mod 26:", len(unit), unit)
print("invers 3 mod 26 =", invers(3, 26), " invers 2 mod 26 =", invers(2, 26))

g0, x0, y0 = egcd(1071, 462)
print("gcd(1071,462) =", g0, " x,y =", x0, y0, " langkah =", langkah(1071, 462))

p, g, a, b = 10007, 5, 1234, 5678
A, B = pow(g, a, p), pow(g, b, p)
print("DH: A =", A, " B =", B, " K =", pow(B, a, p), pow(A, b, p))
print("orde(5) =", orde(5, p), " orde(5^5003) =", orde(pow(5, 5003, p), p))

random.seed(29)
L = [langkah(random.getrandbits(32), random.getrandbits(32))
     for _ in range(20000)]
print("euklid 32-bit: rata %.3f, maks %d, min %d" % (sum(L) / len(L), max(L), min(L)))

p, q, e = 1019, 1031, 3
n, phi = p * q, (p - 1) * (q - 1)
d = invers(e, phi)
print("RSA: n =", n, " bit =", n.bit_length(), " phi =", phi, " d =", d)
for m in (424242, 42):
    c = pow(m, e, n)
    print("  m =", m, " c =", c, " D(c) =", pow(c, d, n), " m^3 =", m ** 3)
kali = pow(1234, e, n) * pow(567, e, n) % n
print("E(1234)E(567) mod n =", kali, " D(kali) =", pow(kali, d, n),
      " 1234*567 mod n =", 1234 * 567 % n)

def tabrakan():
    terlihat, k = {}, 0
    while True:
        m = random.getrandbits(32)
        h = hashlib.sha256(m.to_bytes(4, "big")).digest()[:2]
        k += 1
        if h in terlihat:
            return k
        terlihat[h] = m

T = [tabrakan() for _ in range(20)]
print("tabrakan 16-bit: rata %.1f, min %d, maks %d" % (sum(T) / 20, min(T), max(T)))
print("prediksi sqrt(pi/2 * 2^16) =", round(math.sqrt(math.pi / 2 * 2**16), 1))
\end{verbatim}

Keluarannya, apa adanya:

\begin{verbatim}
unit mod 26: 12 [1, 3, 5, 7, 9, 11, 15, 17, 19, 21, 23, 25]
invers 3 mod 26 = 9  invers 2 mod 26 = None
gcd(1071,462) = 21  x,y = -3 7  langkah = 3
DH: A = 4873  B = 3189  K = 5674 5674
orde(5) = 10006  orde(5^5003) = 2
euklid 32-bit: rata 18.750, maks 32, min 6
RSA: n = 1050589  bit = 21  phi = 1048540  d = 699027
  m = 424242  c = 618000  D(c) = 424242  m^3 = 76355615883580488
  m = 42  c = 74088  D(c) = 42  m^3 = 74088
E(1234)E(567) mod n = 551173  D(kali) = 699678  1234*567 mod n = 699678
tabrakan 16-bit: rata 362.7, min 100, maks 743
prediksi sqrt(pi/2 * 2^16) = 320.8
\end{verbatim}
\end{lab}

Baris pertama menegaskan bagian pertama: dua belas unit modulo $26$, dengan $1$ dan $25$
kebalikan dirinya sendiri dan sepuluh sisanya berpasangan; angka $12$ itu sama dengan
$\varphi(26)$.

Baris kedua dan ketiga memperlihatkan dua wajah Euklid: tiga pembagian untuk $1071$ dan $462$
dengan koefisien Bezout $(-3,7)$, dan rata-rata $18{,}75$ pembagian dengan maksimum $32$ untuk
dua ribu pasangan acak $32$ bit. Baris keempat dan kelima adalah pertukaran kunci yang
dijanjikan: $4873$ dan $3189$ di saluran, $5674$ di kedua ujung, sedangkan $5^{5003}$ berorde
$2$ dan karena itu jebakan bagian lima.

Tiga baris RSA memperlihatkan dua hal berlawanan. Untuk $m=424\,242$ semuanya bekerja: sandi
$618\,000$ kembali menjadi $424\,242$. Untuk $m=42$ dengan $e=3$, sandinya $74\,088$, persis
$42^3$, sehingga tidak ada kunci yang perlu dicari, hanya akar pangkat tiga; baris berikutnya
menunjukkan sifat multiplikatif: hasil kali dua sandi membuka pesan hasil kali.

Empat baris terakhir menguji ramalan ulang tahun: rata-rata $362{,}7$ sampel melawan ramalan
$320{,}8$, dengan sebaran lebar seperti yang diharapkan distribusi eksponensial.

\begin{untukdosen}
Urutan mengajar yang saya sarankan: mulai dari sandi mahasiswa di pembuka, lalu lemma
keterdefinisian dan teorema unit pada pertemuan pertama; sisakan Diffie-Hellman dan RSA untuk
pertemuan kedua, dan taruh Shor serta pasca-kuantum sebagai bacaan pertemuan ketiga. Jangan
mengajarkan Euklid sebelum menanyakan ``kapan pembagian ada''.

Dua miskonsepsi yang hampir selalu muncul. Pertama, mahasiswa mengira $\varphi(n)=n-1$ untuk
semua $n$, lalu kaget ketika $\varphi(26)=12$; sembuhkan dengan meminta mereka mendaftar unit
modulo $26$ di papan sebelum rumus apa pun disebut. Kedua, mereka mengira RSA aman karena $n$
sulit difaktorkan, tanpa menyebut ukuran $n$ dan tanpa menyebut bahwa kesulitan itu asumsi;
minta mereka menuliskan klaim keamanan dalam bentuk ``bila ... maka ...''.

Aktivitas papan tulis: dua kelompok menghitung $3^{12}\bmod26$, satu dengan perkalian berulang
dan satu dengan Fermat kecil, lalu bandingkan waktu dan hasilnya. Tanda kesiapan: mereka bisa
menjelaskan tanpa catatan mengapa $[2]_{26}$ tidak punya kebalikan.
\end{untukdosen}

\section{Ringkasan}
\begin{ringkasan}
\begin{itemize}
    \item Aritmetika modulo $n$ bekerja pada kelas sisa, dan pembagian hanya ada untuk kelas yang
        tidak berbagi faktor dengan $n$; Euklid memberi kebalikannya dalam waktu sebanding
        jumlah digit.
    \item Lemma pengacakan, yaitu perkalian dengan unit hanyalah pengacakan urutan, memberi
        teorema Euler dalam lima baris dan Fermat kecil sebagai akibat. Orde membagi
        $\varphi(n)$, dan akar primitif adalah unsur berorde penuh.
    \item Diffie-Hellman menyepakati kunci lewat kesamaan $g^{ab}$ dari dua arah, dengan keamanan
        bersyarat pada kesulitan logaritma diskret; tanpa autentikasi ia jatuh pada serangan
        orang di tengah.
    \item RSA membalik arah: kunci publik mengunci, kunci rahasia membuka, dan buktinya memakai
        Fermat kecil pada kedua faktor prima. RSA polos deterministik, multiplikatif, dan tidak
        melindungi pesan kecil; bantalan menutup ketiganya.
    \item Fungsi hash memberi sidik jari, MAC memberi keaslian antara dua pihak yang berbagi
        kunci, tanda tangan memberi bukti yang bisa diperiksa siapa saja; ketahanan tabrakan
        hanya separuh panjang keluaran, dan angka $362{,}7$ melawan $320{,}8$ di lab
        menunjukkan ramalan itu bekerja.
    \item Semua jaminan di bab ini bersyarat. Shor meruntuhkan RSA dan seluruh logaritma diskret
        bila komputer kuantum berskala besar ada; sandi simetris dan hash cukup menggandakan
        panjang kunci, dan skema lattice seperti ML-KEM adalah pengganti yang mungkin.
\end{itemize}
\end{ringkasan}

\section{Latihan}
\begin{latihan}
\begin{enumerate}
  \item Nilai klaim ini: ``Karena $2^{100}\equiv1\pmod{101}$, bilangan $101$ terbukti prima.''
        Tentukan benar, salah, atau belum cukup bukti, dan sebutkan apa yang dibuktikan
        kongruensi itu.
  \item Buktikan bahwa hasilkali dua unit modulo $n$ juga unit, lalu simpulkan bahwa $U(n)$
        tertutup terhadap perkalian; tunjukkan juga bahwa bila $a$ bukan unit, ada $b\neq0$
        dengan $ab\equiv0\pmod n$.
  \item Buktikan bahwa untuk prima $p$, $a^2\equiv1\pmod p$ hanya dipenuhi $a\equiv1$ atau
        $a\equiv-1$; jelaskan mengapa hal yang sama gagal untuk modulus $8$, dan sebutkan
        semua akar kuadrat dari satu modulo $8$.
  \item Buktikan $\varphi(pq)=\varphi(p)\varphi(q)$ untuk prima berbeda $p,q$ dengan menghitung
        kelas $0\le a<pq$ yang berbagi faktor dengan $pq$, lalu periksa pada $p=1019$ dan
        $q=1031$.
  \item Hitung $17^{-1}\bmod3120$ dengan Euklid mundur, lalu periksa dengan Python. Ulangi
        untuk $65\,537^{-1}\bmod1\,048\,540$ dan bandingkan jumlah pembagian keduanya.
  \item Hitung $5^{100}\bmod101$, $3^{200}\bmod101$, dan $2^{1000}\bmod101$ memakai Fermat
        kecil, tanpa menghitung pangkat penuh. Periksa ketiganya dengan \texttt{pow(a,k,m)}.
  \item Pada modulus $23$ dengan basis $5$, Alice memakai $a=7$ dan Bob $b=15$. Hitung $A$,
        $B$, dan kunci bersama, lalu hitung kunci yang diperoleh penyadap aktif yang mengganti
        $A$ dengan $11$ di depan Bob, dan jelaskan mengapa dua kunci berbeda muncul.
  \item Jelaskan mengapa RSA dengan $e=3$ dan pesan $m<n^{1/3}$ bisa dibuka tanpa kunci rahasia,
        berapa panjang pesan terpanjang yang terkena pada modulus $1\,050\,589$, dan apa yang
        dilakukan bantalan untuk mencegahnya.
  \item Nilai klaim ini: ``ML-KEM sudah terbukti aman, jadi kita bisa berhenti memakai RSA
        sekarang.'' Sebutkan dua hal yang harus benar agar kalimat pertama masuk akal, dan satu
        alasan mengapa kalimat kedua tetap keliru meski yang pertama benar.
  \item Sebuah sistem memakai hash berkeluaran $48$ bit. Berapa sampel yang dibutuhkan untuk
        peluang tabrakan satu per dua, dan berapa bila keluarannya dipotong menjadi $32$ bit?
\end{enumerate}

Petunjuk untuk tiga butir tersulit. Butir 3: tulis $a^2-1=(a-1)(a+1)$ dan pakai lemma Euklid;
untuk modulus $8$, cari akar dengan mencoba $a=0,\dots,7$. Butir 7: hitung $11^7\bmod23$
sebagai kunci yang dilihat Bob, lalu bandingkan dengan $11^{15}\bmod23$ yang dihitung Alice.
Butir 9: bedakan ``terbukti aman'' dari ``aman terhadap serangan yang diketahui''.
\end{latihan}

\begin{untukpeneliti}
Yang benar-benar terbukti di bab ini adalah bagian matematikanya: teorema Bezout, teorema unit,
teorema Euler, lemma orde, dan teorema RSA semuanya punya bukti lengkap dan tidak bergantung
pada dugaan. Yang hanya dipopulerkan adalah kata ``aman''. Tidak ada teorema yang menyatakan
RSA-2048 sulit; yang ada taksiran biaya algoritme terbaik yang diketahui, dan taksiran itu
berubah setiap kali ada kemajuan kriptanalisis.

Sumber primernya: Diffie dan Hellman (1976), bagian 3--4; Rivest, Shamir, dan Adleman (1978),
bagian 2--4; Shor (1997), bagian 5 untuk reduksi faktorisasi ke pencarian orde; Regev (2005),
bagian 3 untuk reduksi LWE; Bleichenbacher (1998); Gidney dan Ekera (2021); FIPS 203 dan 204.

Batas keberlakuan yang perlu disebut: reduksi keamanan biasanya menghasilkan batas longgar,
model serangan tidak memuat saluran samping, dan serangan padding menunjukkan bahwa selisih satu
pesan galat cukup meruntuhkan skema yang terbukti aman dalam model. Pertanyaan terbuka: seberapa
ketat reduksi untuk ML-KEM, apakah ada algoritme kuantum yang lebih cepat untuk masalah lattice,
dan bagaimana membuktikan keamanan skema yang dipakai bersama data rahasia tanpa membocorkan
pola akses.
\end{untukpeneliti}

\peta{Bab 31 pada Diffie dan Hellman (1976), bagian 3 untuk pertukaran kunci; Rivest, Shamir,
dan Adleman (1978), bagian 2--4 untuk RSA dan tanda tangan; Shor (1997), bagian 5--6 untuk
reduksi faktorisasi ke pencarian orde; Regev (2005), bagian 3 untuk reduksi LWE; Bleichenbacher
(1998); Gidney dan Ekera (2021); FIPS 203 dan 204. Pada Bab 2 untuk bijeksi dan partisi, pada
Bab 3 untuk sifat Archimedes dan bentuk diskret sifat infimum, pada Bab 8 untuk eliminasi
Gauss, pada Bab 26 untuk alasan klaim keamanan tidak bisa berbentuk teorema tanpa syarat, dan
pada Bab 29--30 untuk pola ancaman oracle.}
